% !PRINTED 107
\weberpage{107}
If we take as an example $m = 48$, we have two possibilities, namely:
\[ \begin{aligned}
e_{1,1} &= 1, & e_{1,2} &= 0, & e_{1,3} &= 1\\
e_{1,1} &= 1, & e_{1,2} &= 1, & e_{1,3} &= 1,
\end{aligned} \]
and from these one finds the two groups:
\[ \begin{aligned}
\mathfrak{A}_1 &= 1,\ 17,\ 31,\ 47\\
\mathfrak{A}_2 &= 1,\ \ 7,\ 41,\ 47.
\end{aligned} \]

\begin{center}
\S.\,25.
\end{center}

\begin{center}
\gesperrt{Cubic Abelian Equations.}
\end{center}

We shall conclude these considerations of the cubic and
biquadratic cyclotomic equations with the proof of a remarkable
theorem, which is suited to lend these equations
 a far higher and more general interest,
 and which is a special case of an entirely general theorem
that we shall come to know only later.

The theorem runs thus:

\gesperrt{All cubic and biquadratic Abelian
equations in the field of the rational numbers are
cyclotomic equations.}

The results of the two preceding sections thereby gain
 a higher value. For it then follows that
by the cyclotomic periods $\eta$ formed there \gesperrt{all Abelian
number fields} of the third and fourth degree are represented,
and that fields of this kind other than those discussed there do not
exist.

The possibility of proving this theorem for the special cases of the
cubic and biquadratic equations simply and without
further preparation rests on the fact that in the
fields of the third and fourth roots of unity $R(\varrho)$, $R(i)$ the same
 laws of the decomposition of numbers into their prime factors
hold as in the fields of the rational numbers, as is demonstrated in
\S\S.\,173, 174 of the first volume. It is therefore
also of interest to consider these special cases more closely,
because the course and the goal of the general proof
 to be supplied later are thereby more clearly recognised.
